                                                                     \documentclass[12pt, a4paper]{article}
\usepackage[utf8]{inputenc}
\usepackage{amsmath, amssymb, amsthm, amsfonts}
\usepackage{CJKutf8}
\usepackage{tikz}
\usetikzlibrary{arrows.meta, patterns}
\usepackage[margin=1in]{geometry}
\newtheorem{lemma}{Lemma}

 
\newenvironment{solution}
  {\paragraph{\textit{Solution:}}\par\medskip}
  {\hfill$\blacksquare$\vspace{0.1cm}}

\title{Topics in Geometry and Topology:Assignment 3}
\author{
    \begin{CJK*}{UTF8}{gbsn}
    钟皓宇 (12533556)
    \end{CJK*}
}
\date{\today}

\begin{document}

\maketitle
\section*{Problem 1}
If $\tau$ is a face of a cone $\sigma$, show that the sum of two vectors in $\sigma$ can be in $\tau$ only if both of the summands are in $\tau$. Show conversely that any convex subcone of a cone $\sigma$ satisfying this condition is a face.

\begin{solution}
By definition of face, there exists $u\in \sigma^\vee$ such that $\tau= \sigma\cap u^\perp$. According to the property of face presented on the page 13 of notebook, 
we claim that $\tau^\vee = \sigma^\vee+\mathbb{R}_{\ge 0}\cdot(-u)$. Suppose that $v_1+v_2\in\tau$, equivalently it meaning that 
for any vectors $a\in \sigma^\vee$ and $b\in \mathbb{R}_{\ge 0}\cdot(-u)$, both the values of pairs $\langle v_1+v_2,a \rangle\ge 0$ and $\langle v_1+v_2,b \rangle\ge 0$ are greater than zero.
Because $a\in\sigma^\vee$, then we obtain that $\langle v_i,a\rangle\ge 0$ for $i=1,2$. Hence, we also obtain that 
\begin{align*}
    \langle v_1,b\rangle  \ge \langle v_2,-b\rangle\ge 0
    \text{ and }
    \langle v_2,b\rangle  \ge \langle v_1,-b\rangle\ge 0,
\end{align*}
since $-b\in \mathbb{R}_{\ge 0}\cdot u \subseteq \sigma^\vee$ implying that $v_1,v_2\in \tau$.

Conversely, suppose that $\tau$ is a subcone of $\sigma$ such that the condition holds. 
Let $W=\tau+(-\tau)$ be the linear subspace spanned by $\tau$. We claim that $\tau=\sigma\cap W$. It is clear that $\tau\subseteq \sigma\cap W$ by definition.
Suppose $v\in \sigma\cap W$ then there exists $u_1\in \tau$ and $u_2\in (-\tau)$ such that $v= u_1+u_2$. Since $v-u_2 = u_1\in \tau$, then $v$ as a summand of $u_1$ is contained in $\tau$ and thereofore $\sigma\cap W\subseteq \tau$.
Let $u_0$ be an arbitrary relative interior point of $\tau$ and suppose that $\sigma^\vee$ is generated by $\{m_1,\dots,m_s\}$.
Let 
\begin{align*}
    I = \{ i \mid \langle m_i, u_0 \rangle = 0 \}
\end{align*}
be the index set of $m_i$ annihilating at $u_0$. If $I=\emptyset$, we obtain that $u_0$ is also the relative interior point of $\sigma$. Then for arbitrary $v\in \sigma$ and sufficiently small $\varepsilon>0$, we obtain 
that $u_0 = \varepsilon\cdot v+(u_0-\varepsilon\cdot v)$ is contained in $\tau$. Then $\varepsilon\cdot v$ as a summand of $u_0$ is contained in $\tau$. Hence $v$ is contained in $\tau$ which implies that $\sigma\subseteq \tau$.
Therefore we may assume that $I$ is not empty. Let $m=\sum_{i\in I}m_i$. I claim that $\tau=\sigma\cap m^\perp$.
Suppose that $v$ is an arbitrary element in $\tau$, since $u_0$ is relative interior point, then exists $\varepsilon>0$ such that $u_0-\varepsilon\cdot (v-u_0)\in\tau$. 
Consider the values pairing with $m_i$ for $i\in I$, we obtain that 
\begin{align*}
    -\varepsilon \cdot \langle m_i,v \rangle \ge 0.
\end{align*}
However, $v\in \sigma$ implies that $\langle m_i,v\rangle\ge 0$, and therefore $v\in m^\perp$ meaning that $\tau\subseteq \sigma\cap m^\perp$.
For an arbitrary element $v\in\sigma\cap m^\perp$, there exists sufficiently small $\varepsilon>0$ such that $\langle m_i,u_0-\varepsilon\cdot v\rangle\ge 0$ for $i=1,\dots,s$, equivalently $u_0-\varepsilon\cdot v\in \sigma$.
Since we just need to find $\varepsilon$ such that $\langle m_i,u_0\rangle\ge \langle m_i,\varepsilon\cdot v\rangle$ for $i\notin I$.
Then we consider $\varepsilon\cdot v+(u_0-\varepsilon\cdot v)= u_0\in \tau$, $\varepsilon\cdot v$ as a summand of $u_0$ is contained in $\tau$, which implies that $\sigma\cap m^\perp \subseteq \tau$.
Therefore, our claim $\tau=\sigma\cap m^{\perp}$ implies that $\tau$ is a face of $\sigma$.
\end{solution}



\section*{Problem 2}
Let $\varphi_* : X(\Delta') \to X(\Delta)$ be the map arising from a homomorphism of lattices $\varphi : N' \to N$ mapping $\Delta'$ to $\Delta$, which sends a cone $\sigma'$ of $\Delta'$ into a cone $\sigma$ of $\Delta$. Show that for a cone $\sigma'$ in $\Delta'$, $\varphi_*$ maps the point $x_{\sigma'}$ of $X(\Delta')$ to the point $x_\sigma$ of $X(\Delta)$, where $\sigma$ is the smallest cone of $\Delta$ that contains $\varphi(\sigma')$.
\begin{solution}
    The point $x_{\sigma'}$ can be interpreted as a map
    \begin{align*}
        x_{\sigma'}(m') = \begin{cases} 1, & \text{if } m' \in (\sigma')^\perp \cap M' \\ 0, & \text{otherwise} \end{cases}.
    \end{align*}
    By definition, for arbitrary $m\in M$, we have 
    \begin{align*}
        (\varphi_*(x_{\sigma'}))(m) = x_{\sigma'}(\varphi^*(m))= \begin{cases} 1, & \text{if } \varphi^*(m) \in (\sigma')^\perp \cap M' \\ 0, & \text{otherwise} \end{cases}.
    \end{align*}
    where $\varphi^*:M\to M'$ is the pullback of $\varphi$. Hence, there are equivalent statements
    \begin{align*}
        \varphi^*(m) \in (\sigma')^\perp \cap M' 
        \Longleftrightarrow \langle\varphi^*(m),\sigma'\rangle = 0
        \Longleftrightarrow \langle m,\varphi(\sigma')\rangle = 0
        \Longleftrightarrow m\in \varphi(\sigma')^\perp\cap M.
    \end{align*}
    Let $\sigma$ be the smallest cone of $\Delta$ that contains $\varphi(\sigma')$. I claim that
    \begin{align*}
        \sigma^\vee\cap \varphi(\sigma')^\perp\cap M = \sigma^\perp \cap M.
    \end{align*}
    Since $\varphi(\sigma')\subseteq \sigma $ then $\sigma^\perp \cap M\subseteq \sigma^\vee\cap\varphi(\sigma')^\perp\cap M$ is quite clear. We now prove conversely.
    Suppose that $m\in \sigma^\vee\cap \varphi(\sigma')^\perp\cap M$, we define $\tau=m^\perp \cap \sigma$ as a face of $\sigma$ and therefore $\varphi(\sigma')\subseteq \tau$. For the smallest property, $\sigma$ has to equal $\tau$ 
    since $\tau$ as a face of $\sigma$ is also contained in $\Delta$. Therefore, $m^\perp \cap \sigma = \sigma$, implying $\sigma \subseteq m^\perp$, or equivalently, $m \in \sigma^\perp \cap M$, which implies our claim holds.
    For $m\in \sigma^\vee$, we find that 
    \begin{align*}
        (\varphi_*(x_{\sigma'}))(m) = x_{\sigma}(m).
    \end{align*}
    Since $\varphi_*(x_{\sigma'})$ and $x_{\sigma}$ agree on $\sigma^\vee$ then they represente the same closed point in $U_\sigma$.
\end{solution}


\section*{Problem 3}
Show that the distinguished point $x_\tau$ is the origin (identity) of the torus $O_\tau$. Show that $O_\tau$ is the unique closed $T_N$-orbit in $U_\tau$.
\begin{solution}
    For arbitrary $x\in O_{\tau}$, by definition, we interprete it as a group homomorphism
    \begin{align*}
        x:\tau^\perp\cap M\to \mathbb{C}^*.
    \end{align*}
    Then by definition we obtain that 
    \begin{align*}
        x_{\tau}(m)= 
\begin{cases} 
1, & \text{if } m \in \tau^\perp \\ 
0, & \text{if } m \notin \tau^\perp 
\end{cases}
    \end{align*}
    and therefore $x_{\tau}(m)\cdot x(m)= x(m)$ implies $x_\tau$ is the identity of torus $O_\tau$.
    Here we use a method of extending by zero to redefine $x$. Explicitly, 
    \begin{align*}
        \tilde{x}(m)= 
        \begin{cases} 
x(m), & \text{if } m \in \tau^\perp \\ 
0, & \text{if } m \notin \tau^\perp 
\end{cases}
    \end{align*}
    which is well-defined as we have discussed in our class. Therefore, we should not be confused by the difference in domains of $x$ and $x_\tau$.

    To prove $O_{\tau}$ is the unique closed $T_N$-orbit in $U_\tau$,
    consider the orbit decomposition 
    \begin{align*}
        U_\tau = \coprod_{\sigma\prec\tau} O_{\tau}
    \end{align*}
    Then every $T_N$-orbit in $U_\tau$ must be of the form $O_{\sigma}$ where $\sigma$ is a face of $\tau$.
    Suppose $O_\sigma$ is closed, then we obtain that 
    \begin{align*}
        O_\sigma = \overline{O}_\sigma = V(\sigma) \cap U_{\tau} =  \coprod_{\sigma\prec\gamma\prec\tau} O_{\gamma}
    \end{align*}
    where $\overline{O}_\sigma$ denotes the closure in $U_\tau$. And therefore we obtain that $\sigma=\tau$, equivalently, $O_\tau$ is the unique closed $T_N$-orbit in $U_\tau$.
\end{solution}

\end{document}